\documentclass[11pt,a4paper]{article}
\usepackage[margin=2.2cm]{geometry}
\usepackage{amsmath,amssymb,slashed}
\usepackage{booktabs,array}
\usepackage{xcolor}
\usepackage[colorlinks=true,linkcolor=blue,urlcolor=blue,citecolor=blue]{hyperref}

\newcommand{\ep}{\epsilon}
\newcommand{\dd}{\mathrm{d}}
\newcommand{\kpar}{k_{\parallel}}
\newcommand{\kperp}{k_{\perp}}
\newcommand{\note}[1]{\par\smallskip\noindent\colorbox{gray!12}{\parbox{\dimexpr\linewidth-2\fboxsep}{\small #1}}\par\smallskip}

\title{One-loop $\times$ one-loop interferences\\[2pt]\large factorisable two-loop families, the $k_1\cdot k_2$ numerators,
and complex conjugation}
\author{DY\_P2B --- notes\_Larin/}
\date{\today}

\begin{document}
\maketitle

\section*{Summary}

In the real--virtual--virtual contribution at N$^3$LO, the product of two one-loop amplitudes,
$\mathcal{A}^{(1)}\,\mathcal{A}^{(1)*}$, contains two loop momenta: $k_1$ in the amplitude and $k_2$ in the
complex-conjugate amplitude. No propagator depends on both, but the spin and polarisation sums contract Lorentz
indices of the two loops, so the numerators contain $k_1\cdot k_2$. The integrals are therefore not plain products
of one-loop integrals, but they belong to a \emph{factorisable} two-loop family. We show:
\begin{enumerate}
\item each family has 12 propagators: 3 cut, 4 in $k_1$, 4 in $k_2$, and one irreducible numerator for
$k_1\cdot k_2$ (section~\ref{sec:count});
\item every integral with powers of $k_1\cdot k_2$ is a combination of products of one-loop integrals without it
(Lorentz covariance, section~\ref{sec:k1k2}); the reduction of the two-loop family therefore yields masters that are
products $M_a(k_1)\,M_b(k_2)$;
\item the reduction never mixes $k_1$ and $k_2$ propagators, and its coefficients are real, so the opposite $i0$
prescription of the conjugate loop is applied at the end, $M_a\,M_b^*$; symmetries exchanging $k_1$ and $k_2$ must
not be used (section~\ref{sec:conj}).
\end{enumerate}
No tensor reduction is needed: the reduction of the two-loop family does it.

\section{Setting}

One emission: $q(p_1)+\bar q(p_2)\to V(q)+g(r_1)$ and crossings, $s=(p_1+p_2)^2$, $q^2=m^2=zs$, and the rapidity
fixed by $u$. The cut conditions fix $r_1$ completely,
\begin{equation}
2\,p_1\cdot r_1=s\,\frac{1-uz}{1+u},\qquad 2\,p_2\cdot r_1=s\,\frac{u-z}{1+u},
\end{equation}
so every invariant is a function of $(s,z,u)$. The interference
\begin{equation}
\sum_{\rm spins,\,pols}\mathcal{A}^{(1)}(k_1)\;\mathcal{A}^{(1)}(k_2)^*
\end{equation}
is written with the loop momentum $k_1$ in the amplitude and $k_2$ in the conjugate. The amplitude loop has Feynman
propagators $1/(\ell^2+i0)$; the conjugate loop has $1/(\ell^2-i0)$.

\paragraph{Where $k_1\cdot k_2$ comes from.} The open quark line of the amplitude and that of the conjugate join into
one Dirac trace, e.g.\ $\mathrm{Tr}[\,\slashed{p}_2\,\Gamma(k_1)\,\slashed{p}_1\,\bar\Gamma(k_2)\,]$, and the gluon
polarisation sum $-g^{\mu\nu}$ and the leptonic tensor contract a Lorentz index of the amplitude with one of the
conjugate. A $\slashed{k}_1$ from $\Gamma$ and a $\slashed{k}_2$ from $\bar\Gamma$ in the same trace give $k_1\cdot k_2$;
two pairs give $(k_1\cdot k_2)^2$, and so on.

\section{Counting: why 12 propagators}\label{sec:count}

Integration momenta $\{r_1,k_1,k_2\}$ (the emission and the two loops), external momenta $\{p_1,p_2\}$.
The scalar products that involve an integration momentum are
\begin{center}
\begin{tabular}{lll}
\toprule
group & scalar products & number \\
\midrule
emission & $r_1^2,\ r_1\cdot p_1,\ r_1\cdot p_2$ & 3 \\
loop 1 & $k_1^2,\ k_1\cdot p_1,\ k_1\cdot p_2,\ k_1\cdot r_1$ & 4 \\
loop 2 & $k_2^2,\ k_2\cdot p_1,\ k_2\cdot p_2,\ k_2\cdot r_1$ & 4 \\
mixed & $k_1\cdot k_2$ & 1 \\
\midrule
 & & 12 \\
\bottomrule
\end{tabular}
\end{center}
A family needs 12 linearly independent denominators (or numerators) to express all of them:
\begin{itemize}
\item 3 cut propagators: $r_1^2$, $(p_1+p_2-r_1)^2-m^2$, and the linear rapidity denominator;
\item 4 propagators of the box in $k_1$: $k_1^2,\ (k_1+P_1)^2,\ (k_1+P_2)^2,\ (k_1+P_3)^2$, with $P_i$ the partial sums
of the external momenta in one cyclic ordering of the legs $p_1,p_2,r_1,q$; they are independent in
$k_1^2,k_1\cdot p_1,k_1\cdot p_2,k_1\cdot r_1$;
\item 4 propagators of the box in $k_2$, the same for $k_2$;
\item 1 irreducible numerator for $k_1\cdot k_2$, e.g.\ $(k_1+k_2)^2$: it never appears in a denominator.
\end{itemize}
$3+4+4+1=12$. Since the cuts fix $r_1$, one may equally reduce in $k_1,k_2$ only, at fixed $(s,z,u)$: 9 scalar
products, $4+4$ propagators and the same numerator. The two descriptions have the same loop part.

\note{One loop $\times$ tree (RV) has no such numerator: the conjugate has no loop momentum, and the contraction only
produces $k_1\cdot p_i$ and $k_1\cdot r_1$, which the box propagators express.}

\section{Factorisable families}\label{sec:fact}

A two-loop family is \emph{factorisable} if its propagators split into two groups depending on disjoint loop
momenta: here the $k_1$-box and the $k_2$-box. As a graph it is two one-loop graphs with no propagator in common.
An integral of the family is
\begin{equation}
I(\vec a;\vec b;n)=\int\!\frac{\dd^dk_1}{i\pi^{d/2}}\frac{\dd^dk_2}{i\pi^{d/2}}\;
\frac{(k_1\cdot k_2)^n}{\prod_{j}D_j(k_1)^{a_j}\;\prod_{j}E_j(k_2)^{b_j}},
\end{equation}
($D_j$ the $k_1$-box propagators, $E_j$ the $k_2$-box ones, negative $a_j,b_j$ allowed). For $n=0$ it is a product
of one-loop integrals,
\begin{equation}
I(\vec a;\vec b;0)=J_1(\vec a)\;J_2(\vec b),\qquad
J_1(\vec a)=\int\!\frac{\dd^dk_1}{i\pi^{d/2}}\frac{1}{\prod_jD_j(k_1)^{a_j}},
\end{equation}
because the integrand is a product of a function of $k_1$ and a function of $k_2$. The only coupling between the two
loops is the numerator $(k_1\cdot k_2)^n$.

\section{Removing $k_1\cdot k_2$}\label{sec:k1k2}

\subsection{Lorentz covariance}

Write the $k_1$-part of the integrand as $f(k_1)=1/\prod_jD_j(k_1)^{a_j}$. It depends on $k_1$ only through $k_1^2$ and
$k_1\cdot p_i$, with $\{p_i\}=\{p_1,p_2,r_1\}$ (three linearly independent external momenta). Split $k_1$ into its
component in the space spanned by the $p_i$ and the rest:
\begin{equation}
k_1^\mu=\kpar^\mu+\kperp^\mu,\qquad
\kpar^\mu=\sum_{i,j}p_i^\mu\,(G^{-1})_{ij}\,(k_1\cdot p_j),\qquad G_{ij}=p_i\cdot p_j ,
\end{equation}
and the same for $k_2$. ``Parallel'' means: the projection of $k_1$ onto the three-dimensional subspace spanned by
$p_1,p_2,r_1$; it is the combination of the $p_i$ with the same scalar products as $k_1$, since contracting with
$p_l$ gives $\kpar\cdot p_l=\sum_{ij}G_{li}(G^{-1})_{ij}(k_1\cdot p_j)=k_1\cdot p_l$. ``Perpendicular'' means:
$\kperp=k_1-\kpar$ is orthogonal to every $p_l$, $\kperp\cdot p_l=0$.
$\kpar$ and $\kperp^2=k_1^2-\kpar^2$ are functions of the scalar products of the family,
and $\kperp$ lives in the $(d-3)$-dimensional space orthogonal to all $p_i$. The projector onto that space is
\begin{equation}
g_\perp^{\mu\nu}=g^{\mu\nu}-\sum_{i,j}p_i^\mu\,(G^{-1})_{ij}\,p_j^\nu ,\qquad
g_\perp^{\mu\nu}p_{l\,\nu}=0,\qquad g_\perp^{\mu\nu}g_{\perp\,\mu\nu}=g_\perp{}^\mu{}_\mu=d-3 .
\end{equation}
We need the tensor integrals
\begin{equation}
T^\mu=\int\!d^dk_1\,f(k_1)\,\kperp^\mu ,\qquad T^{\mu\nu}=\int\!d^dk_1\,f(k_1)\,\kperp^\mu\kperp^\nu .
\end{equation}

\paragraph{General form.} The integrand of $T^{\mu\nu}$ depends on $k_1$ only through scalar products, and the
integration runs over all of $k_1$, so $T^{\mu\nu}$ is a Lorentz-covariant function of the $p_i$:
$\Lambda^\mu{}_\rho\Lambda^\nu{}_\sigma T^{\rho\sigma}(p)=T^{\mu\nu}(\Lambda p)$. A covariant tensor built from vectors
$p_i$ is a combination of products of the metric, of the $p_i$, and of the Levi-Civita tensor, with coefficients that
are functions of the invariants $p_i\cdot p_j$. The Levi-Civita tensor does not enter, for two independent reasons:
\begin{enumerate}
\item $T^{\mu\nu}$ is symmetric by construction, while every rank-two structure with one Levi-Civita tensor, such as
$\epsilon^{\mu\nu\rho\sigma}p_{1\rho}p_{2\sigma}$, is antisymmetric in $\mu\nu$;
\item $f$ is parity even. A reflection of one direction orthogonal to all $p_i$ (it exists for $d>3$) leaves the
$p_i$ and $f$ unchanged, so $R^\mu{}_\rho R^\nu{}_\sigma T^{\rho\sigma}=T^{\mu\nu}$, while it flips the sign of every
structure with one Levi-Civita tensor ($\det R=-1$). This excludes such structures at every rank, $T^\mu$ included.
\end{enumerate}
Hence
\begin{equation}
T^\mu=\sum_iC_i\,p_i^\mu ,\qquad
T^{\mu\nu}=A\,g^{\mu\nu}+\sum_{i,j}B_{ij}\,p_i^\mu p_j^\nu ,\qquad B_{ij}=B_{ji}.
\end{equation}

\paragraph{Transversality.} $\kperp\cdot p_l=0$ holds before integration, so $T^\mu p_{l\,\mu}=0$ and
$T^{\mu\nu}p_{l\,\nu}=0$ for every $l$. For the vector, $\sum_iC_iG_{il}=0$ for all $l$, and $G$ is invertible, so
$C_i=0$. For the tensor,
\begin{equation}
A\,p_l^\mu+\sum_{i,j}B_{ij}\,G_{jl}\,p_i^\mu=0
\quad\Longrightarrow\quad A\,\delta_{il}+(BG)_{il}=0
\quad\Longrightarrow\quad B=-A\,G^{-1},
\end{equation}
since the $p_i$ are linearly independent. Hence $T^{\mu\nu}=A\,g_\perp^{\mu\nu}$: $g_\perp$ is the only combination of
$g$ and the $p_i$ orthogonal to all $p_i$.

\paragraph{Trace.} Contracting with $g_{\mu\nu}$ and using $g_\perp{}^\mu{}_\mu=d-3$,
$A=\frac{1}{d-3}\int d^dk_1\,f(k_1)\,\kperp^2$, an integral of the family since $\kperp^2=k_1^2-\kpar^2$ is a
function of its scalar products. Hence, inside the integral,
\begin{equation}
\kperp^\mu\to0,\qquad
\kperp^\mu\kperp^\nu\to\frac{\kperp^2}{d-3}\,g_\perp^{\mu\nu}.
\end{equation}
The same steps at higher rank give zero for every odd product of $\kperp$'s and symmetrised products of $g_\perp$ for
every even one, e.g.\ $\kperp^\mu\kperp^\nu\kperp^\rho\kperp^\sigma\to\frac{(\kperp^2)^2}{(d-3)(d-1)}
\big(g_\perp^{\mu\nu}g_\perp^{\rho\sigma}+g_\perp^{\mu\rho}g_\perp^{\nu\sigma}+g_\perp^{\mu\sigma}g_\perp^{\nu\rho}\big)$.

\paragraph{The same result without the general form.} At fixed $\kpar$, i.e.\ at fixed $k_1\cdot p_i$, the integrand
depends on $\kperp$ only through $\kperp^2=k_1^2-\kpar^2$. The integral over $\kperp$ is therefore invariant under the
rotations of the $(d-3)$-dimensional orthogonal space, and $T^{\mu\nu}$, which lives entirely in that space, must be a
symmetric invariant tensor of it: a multiple of its metric $g_\perp$.

\paragraph{Levi-Civita tensors in the numerator.} The traces with $\gamma_5$ can leave Levi-Civita tensors contracted
with $k_1$, e.g.\ $\epsilon(k_1,p_1,p_2,r_1)=\epsilon_\mu(p_1,p_2,r_1)\,\kperp^\mu$. They are not part of $f$: they are
external tensors contracted with $\kperp^\mu$ or $\kperp^\mu\kperp^\nu$, and the replacements above apply to them as
to $k_2$. The derivation needs only $f$, the denominators and the scalar-product numerators, to be parity even, which
always holds.

\subsection{Rank one}

Expanding,
\begin{equation}
k_1\cdot k_2=\kpar\cdot k_{2\parallel}+\kpar\cdot k_{2\perp}+\kperp\cdot k_{2\parallel}+\kperp\cdot k_{2\perp}
=\kpar\cdot k_{2\parallel}+\kperp\cdot k_{2\perp},
\end{equation}
because the cross terms vanish identically: $\kpar$ is a combination of the $p_i$ and $k_{2\perp}$ is orthogonal to
all of them (and vice versa). The first term is
$\kpar\cdot k_{2\parallel}=\sum_{ij}(k_1\cdot p_i)(G^{-1})_{ij}(k_2\cdot p_j)$; the second is odd in $\kperp$ and drops
out:
\begin{equation}
I(\vec a;\vec b;1)=\sum_{i,j}(G^{-1})_{ij}\;
\Big[\int f(k_1)\,(k_1\cdot p_i)\Big]\;\Big[\int g(k_2)\,(k_2\cdot p_j)\Big].
\end{equation}
Each bracket is a one-loop integral of the family without $k_1\cdot k_2$: $k_1\cdot p_i$ is a combination of box
propagators and constants. So $I(\vec a;\vec b;1)$ is a sum of products $J_1\,J_2$ with rational coefficients.

\subsection{Rank two}

$(k_1\cdot k_2)^2=(\kpar\cdot k_{2\parallel})^2+2(\kpar\cdot k_{2\parallel})(\kperp\cdot k_{2\perp})
+(\kperp\cdot k_{2\perp})^2$. The middle term is odd in $\kperp$; the last one gives
$\kperp^\mu\kperp^\nu k_{2\perp\mu}k_{2\perp\nu}\to\kperp^2\,k_{2\perp}^2/(d-3)$. Hence
\begin{equation}
(k_1\cdot k_2)^2\;\to\;(\kpar\cdot k_{2\parallel})^2+\frac{1}{d-3}\,\kperp^2\,k_{2\perp}^2 ,
\end{equation}
and both terms are products of a function of the $k_1$ scalar products and one of the $k_2$ scalar products.
Any power $n$ works the same way.

\subsection{Can $k_1\cdot k_2$ survive in a master?}

No relation is missing: the identities above are Lorentz-invariance identities, and these are linear combinations of
the integration-by-parts identities of the family \cite{Lee}. A complete reduction of the two-loop family therefore
contains them: every integral with $n>0$ reduces to integrals with $n=0$, which are products. Which integrals are
called masters is a choice of basis: an ordering that ranks numerator powers as more complicated (the usual one)
never picks an integral with $n>0$. The masters are then products,
\begin{equation}
\{\text{masters}\}=\{M_a(k_1)\,M_b(k_2)\},
\end{equation}
with $\{M_a\}$ the one-loop masters of the $k_1$-box and $\{M_b\}$ those of the $k_2$-box: the IBP identities
generated by $\partial/\partial k_1\cdot v$ with $v\in\{k_1,p_i\}$ act on the $k_1$ factor only and are exactly the
one-loop IBPs, so the master basis of the product is the product of the two bases.

\paragraph{Sanity checks.}
\begin{itemize}
\item The coefficients contain $1/\det G$ and $1/(d-3)$. With $p_1^2=p_2^2=r_1^2=0$,
\begin{equation}
\det G=2\,(p_1\cdot p_2)(p_1\cdot r_1)(p_2\cdot r_1)=\frac{s^3}{4}\,\frac{(1-uz)(u-z)}{(1+u)^2},
\end{equation}
non-zero inside the physical region and vanishing on the collinear boundaries $u=z$ and $u=1/z$, where $r_1$ is
parallel to $p_2$ or $p_1$ and the three momenta no longer span three directions. These are the boundaries where the
expansion in distributions is needed anyway. $d-3=1-2\ep$ is harmless.
\item In a subsector the $k_1$-integral may depend on fewer external momenta (a triangle or a bubble); the same
argument holds with the Gram matrix of those momenta.
\item A direct check, on a family with both boxes in the ordering $p_1,p_2,r_1$: eight integrals, with $n=0,1,2$,
one dot, and in subsectors, computed in two ways, (i) with the rank-one and rank-two formulas above followed by the
one-loop reduction of each factor, and (ii) by the reduction of the two-loop family without $k_1\leftrightarrow k_2$
symmetries (section~\ref{sec:conj}). The two agree term by term as ordered products $M_a(k_1)\,M_b(k_2)$, and no master
carries a power of $k_1\cdot k_2$.
\end{itemize}

\section{Complex conjugation}\label{sec:conj}

The $k_1$ propagators carry $+i0$, the $k_2$ propagators $-i0$. This is consistent with the reduction:
\begin{enumerate}
\item The IBP identities are algebraic: they do not depend on the $i0$ prescription, and their coefficients are
rational functions of $s,z,u,d$ with real coefficients.
\item No identity produces a propagator depending on both $k_1$ and $k_2$: $\partial/\partial k_1$ acts only on the
$k_1$ propagators and on the numerator $k_1\cdot k_2$. Every term of the reduction keeps the split (normal $k_1$
factor) $\times$ (conjugate $k_2$ factor).
\end{enumerate}
Hence one reduces as if both prescriptions were $+i0$, and evaluates each master as
\begin{equation}
M_a(k_1;+i0)\;M_b(k_2;-i0)=M_a\;M_b^{*},
\end{equation}
with $M_b$ the usual (Feynman) one-loop master at the same real kinematics.

\paragraph{Symmetries.} Symmetries acting on $k_1$ alone (shifts and reflections of the $k_1$-box) or on $k_2$ alone
preserve the split and may be used. A symmetry that exchanges $k_1$ and $k_2$ (possible when both boxes have the same
ordering) must not be used: it would identify $M_a\,M_b^*$ with $M_b\,M_a^*$, i.e.\ move the conjugation to the
other factor. The exchange relates instead the two interferences with amplitude and conjugate swapped,
$\mathcal{A}_i^{(1)}\mathcal{A}_j^{(1)*}=\big(\mathcal{A}_j^{(1)}\mathcal{A}_i^{(1)*}\big)^*$, which is used when the
contributions are summed, not inside the reduction.

An automatic search for sector relations finds this exchange: for two boxes of the same ordering it maps, for
instance, (box in $k_1$)$\times$(bubble in $k_2$) onto (bubble in $k_1$)$\times$(box in $k_2$), so the reduced
integral is correct as a number for $+i0$ in both loops but no longer correct once the conjugation is applied to the
$k_2$ factor. The reduction of the family is therefore run with sector relations and sector symmetries switched off;
every sector is then reduced on its own, and the masters are the $4\times4=16$ ordered products. In the check of
section~\ref{sec:k1k2} the reduction with the exchange allowed agrees with the formulas of section~\ref{sec:k1k2} for
all eight integrals only after symmetrising each product, $M_a M_b\to\tfrac12(M_aM_b+M_bM_a)$, and as ordered products
for one of them; without the exchange all eight agree as ordered products.

\section{Families and reductions}

\begin{itemize}
\item Families: (ordering of the $k_1$-box) $\times$ (ordering of the $k_2$-box), 3 orderings each (legs
$p_1,p_2,r_1$ around the box, $q$ fixed, up to reflection), at most 9 families of 12 propagators.
\item Two different reductions enter the single-emission virtual corrections: the one-loop boxes (one loop $\times$
tree, NNLO and N$^3$LO) and the factorisable two-loop families of this note (one loop $\times$ one loop).
\item Masters: products of one-loop masters, each needed to the order in $\ep$ required by the product, with the
conjugation applied to the $k_2$ factor. With massless internal lines the triangles reduce to bubbles, so each box
has four masters: the box and the three bubbles in the invariants of its two-particle channels and in $q^2$ (for the
ordering $p_1,p_2,r_1$: $(p_1+p_2)^2$, $(p_2+r_1)^2$, $q^2$); a family has $4\times4=16$ masters.
\end{itemize}

\begin{thebibliography}{9}
\bibitem{Lee} R.~N.~Lee, \emph{Group structure of the integration-by-part identities and its application to the
reduction of multiloop integrals}, JHEP 07 (2008) 031, \href{https://arxiv.org/abs/0804.3008}{arXiv:0804.3008}.
\end{thebibliography}

\end{document}
